\section{Introduction}\label{sec:introduction}

Let $G$ be a connected reductive group over $\C$, let $B$ be a Borel subgroup, and let $P\supseteq B$ be a parabolic subgroup. A Schubert variety $X\subseteq G/P$ is normal and Cohen--Macaulay \cite{Ram85,Bri05}. Its canonical sheaf is therefore a rank-one reflexive sheaf, and its Gorenstein locus is the locus on which that sheaf is invertible \cite[Chapter~3]{BH93}. We ask whether this local obstruction can be detected at the generic points of the maximal singular components.

Write $\M(X)=\Irr(\Sing X)$, and denote the generic point of $Z\in\M(X)$ by $\eta_Z$. Define
\begin{equation}\label{eq:bad}
 \Bbad(X)=\bigcup_{\substack{Z\in\M(X)\\ \OO_{X,\eta_Z}\text{ is not Gorenstein}}} Z.
\end{equation}
This definition concerns the ambient local ring of $X$, rather than the intrinsic singularities of $Z$. Openness of the Gorenstein locus gives $\Bbad(X)\subseteq\NGor(X)$.

Woo and Yong classified globally Gorenstein Schubert varieties in type~A \cite{WY06}. Their Corollary~1 states that a type~A Schubert variety is globally Gorenstein if and only if its local rings at the generic points of all maximal singular components are Gorenstein. Their Conjecture~1 asks for the following pointwise strengthening.

\begin{conjecture}[Woo--Yong]\label{conj:wy}
If $X$ is a Schubert variety in a type~A flag variety, then $\NGor(X)=\Bbad(X)$.
\end{conjecture}

All equalities of loci in this paper are set-theoretic and include nonclosed points; they do not assert equality of scheme structures. Our convention is $X_w=\overline{BwB/B}$ and $\dim X_w=\ell(w)$; some original references use the opposite indexing convention.

The distinction between the global classification and Conjecture~\ref{conj:wy} is substantial. Removing some singular components usually leaves an open subset that is not a Schubert variety, so the global classification cannot simply be reapplied. The issue includes varieties with both generically Gorenstein and generically non-Gorenstein maximal singular components. A proof must handle points on the former components that lie on none of the latter.

\subsection{Divisor detection in type A}
For every scheme point $x\in X$, localization gives
\begin{equation}\label{eq:rho}
 \rho_x:\Cl(\OO_{X,x})\longrightarrow
 \bigoplus_{\substack{Z\in\M(X)\\x\in Z}}\Cl(\OO_{X,\eta_Z}).
\end{equation}
The condition $x\in Z$ makes $\eta_Z$ a generalization of $x$, so each component map is defined. An empty direct sum is the zero group.

\begin{theorem}[Type A divisor detection]\label{thm:mainA}
Let $X$ be a Schubert variety in a type~A complete or partial flag variety. For every scheme point $x\in X$, the map $\rho_x$ is injective. Equivalently, every Weil divisor $D$ on $X$ satisfies
\begin{equation}\label{eq:noncartier}
 \{x:D\text{ is not Cartier at }x\}
 =\bigcup_{\substack{Z\in\M(X)\\ D\text{ is not Cartier at }\eta_Z}} Z.
\end{equation}
\end{theorem}

Taking $D=K_X$ proves Conjecture~\ref{conj:wy}. Allowing all Weil divisors also gives the geometric nonfactorial-locus assertion of \cite[Conjecture~6.14]{WY08} and \cite[Conjecture~7.39]{WY23}. These consequences use the integral statement, including classes of infinite order.

The proof in Section~\ref{sec:type_a} inverts a coordinate associated with the first difference between two permutations. The resulting open subset is an actual product of a smaller Kazhdan--Lusztig slice and $\Gm$, and its complementary divisor is reduced.

Maximal elements of a filtered Bruhat intersection give rank-one determinantal slices in which two specified vertical components are the two rulings through one matrix entry. They detect successive integer differences of divisor coefficients. A separate prefix argument connects a horizontal component to the same comparison. Nagata localization and induction then prove injectivity. No cancellation of affine factors or identification of a tangent cone with the original slice is used.

\subsection{Finite-order detection in arbitrary type}
The same conclusion for all classes fails in other root systems. There is, however, a uniform assertion for finite-order classes.

\begin{theorem}[Torsion detection]\label{thm:maintorsion}
Let $X$ be a Schubert variety in $G/P$ for a finite crystallographic root system, possibly reducible. For every scheme point $x\in X$, the restriction of $\rho_x$ to $\Cl(\OO_{X,x})_{\tors}$ is injective.
\end{theorem}

The torsion hypothesis is imposed on the source class. A calculation after tensoring with $\Q$ would not establish this assertion.

\begin{corollary}[Indices and canonical loci]\label{cor:indices}
Suppose that $D$ is $\Q$-Cartier at $x$. Then
\[
 \ind_x(D)=\mathrm{lcm}_{Z\ni x}\ind_{\eta_Z}(D),
\]
where $Z$ ranges over $\M(X)$ and the empty least common multiple is $1$. On the $\Q$-Gorenstein locus, $X$ is Gorenstein at $x$ if and only if it is Gorenstein at every $\eta_Z$ with $x\in Z$. On all of $X$,
\[
 \NGor(X)=\NonQGor(X)\cup\Bbad(X).
\]
\end{corollary}

The proof of Theorem~\ref{thm:maintorsion} has a different geometric mechanism. In dimension at least four, finite covers of a strict local puncture are detected on its coatom boundary. The topology of an open Bruhat interval supplies compatible sheet labels; a local Lefschetz theorem lifts the resulting splitting. A second argument bounds torsion by degrees of actual integral cycles in a Bott--Samelson fibre.

Low-dimensional slices reduce to Levi subgroups of bounded rank. The simply laced and type~B cases follow from these arguments and low-rank bases. In type~C, the surviving strict torsion models are scalar second-Veronese cones, which have isolated singularities. Type~$F_4$ requires a finite certificate and type~$G_2$ an alternating-word argument. Throughout, field extensions of $\C$ are retained because a Schubert cell's generic point has a nonconstant residue field.

\subsection{The hypothesis beyond type A}
\begin{theorem}[A type D family]\label{thm:mainD}
For every $r\ge k\ge3$, there is a Schubert variety $X\subset\OG(k,2r+4)$ of dimension $k(r-k+2)$ whose singular locus is irreducible of codimension three and generically an ordinary threefold node. Nevertheless,
\[
 \NGor(X)=\NonQGor(X)\simeq\Gr(k,r),
 \qquad \codim_X\NGor(X)=2k.
\]
The same phenomenon occurs after passage to complete flags.
\end{theorem}

In this family $\Bbad(X)=\varnothing$. The smallest member has type~$D_5$, dimension six, and one non-Gorenstein point; the member $(r,k)=(5,4)$ has type~$D_7$, dimension twelve, and non-Gorenstein locus $\PP^4$. They contradict the unrestricted general-type formulation in \cite[Conjecture~0.1]{Per09}, while satisfying the scopes of the type~A and finite-order results: the obstructing canonical class has infinite order.

\subsection{Relation to earlier work}
Perrin proved the locus statement for minuscule Schubert varieties, including ordinary Grassmannian Schubert varieties \cite[Theorem~0.3 and Corollary~0.4]{Per09}. Woo--Yong's interval-pattern framework describes local properties through Bruhat intervals \cite{WY08}; their survey distinguishes the pointwise conjecture from the global classification and records verification through $S_6$ \cite[Conjecture~7.31 and Exercise~7.32]{WY23}.

The product charts used here come from \cite{WY08,KWY13}. Maximal singular components in type~A were classified independently in \cite{BW03,Cor03,KLR03,Man01b}; their generic local models are described in \cite{Cor03,Man01a}. Our filtered construction supplies the additional marked divisors required for coefficient comparisons.

Combining the local Grassmannian reduction for vexillary opposite Schubert varieties in \cite[Theorem~7.1 and Section~5]{AIJK23} with Perrin's theorem gives the corresponding covexillary locus result. Da Silva studies boundary blow-ups and gives a Demazure-subword criterion \cite[Theorem~0.2 and Corollary~0.7]{DS19}.

Li--Rietsch--Yang study global Cartier and $\Q$-Cartier classes on $G/P$ Schubert varieties \cite{LRY25}, addressing the global problem in \cite[Problem~8.24]{WY23}. Kimura determines canonical trace ideals and non-Gorenstein loci for Grassmannian Schubert cycles and determinantal rings \cite{Kim26}. These results have related but different scopes; none states the local injection \eqref{eq:rho} for all ordinary local rings of an arbitrary full-flag Schubert variety.

PermPAL \cite{PermPAL}, based on the framework of \cite{ABC26}, provides specifications and enumerations of permutation avoidance classes. Its class $\operatorname{Av}(3412,52341,635241)$ occurs in the iterated Grassmannian-bundle characterization of Alland--Richmond \cite{AR18}. This observation supplies combinatorial context and is not an input to either detection theorem.

The proof dependencies are explicit. In type~A, Theorem~\ref{ta:filtered} identifies the marked ruling divisors, Proposition~\ref{ta:horizontal} handles the horizontal component, and Lemma~\ref{ta:differences} and Theorem~\ref{ta:affine} complete the coefficient argument.

In arbitrary type, the Levi reduction, curve-degree construction, and boundary theorem provide the geometric inputs, while the type~C argument supplies the unbounded induction. The type~D construction is independent of these steps. Section~\ref{sec:reproducibility} and the accompanying source package specify the finite computations and their reconstruction.
